% GENERATED FILE. Edit canonical lessons, then run npm run generate:books.
% canonical-source-hash: fnv1a64:1d97033cdbf4f3cc
% canonical-lessons: LA-C461-thermopolium, LA-C461-calda, LA-C461-poto, LA-C461-fumosus, LA-C461-immundus

\chapter{Hot-drinks bar, Hot water, To drink, Smoky, Filthy}
\label{ch:la-hot-drinks-bar-hot-water-to-drink-smoky-filthy}
\hlchaptermodality{461}

\begin{chapteropening}
\textbf{By the end of this chapter:} \emph{I can talk about eating and drinking in town: a hot-drinks bar, hot water for wine, drinking, a smoky room and a filthy table.}

Complete the last lesson of chapter 461: I can talk about eating and drinking in town: a hot-drinks bar, hot water for wine, drinking, a smoky room and a filthy table.
\end{chapteropening}

\section[thermopōlium, thermopōliī]{thermopōlium, thermopōliī — a hot-drinks bar, a tavern selling warm drinks}

\label{lesson:LA-C461-thermopolium}



\begin{quote}
Before the new one: say the Latin for to go back and forth, then the Latin for near the city.
\end{quote}

\subsection*{You'll want to know: thermopōlium, thermopōliī}
\textbf{thermopōlium, thermopōliī} — \textquotedblleft{}a hot-drinks bar, a tavern selling warm drinks\textquotedblright{}.

\begin{quote}
Fabrī in thermopōliō vīnum calidum bibunt.
\end{quote}

\begin{quote}
\emph{The workmen drink warm wine in the hot-drinks bar.}
\end{quote}

Second declension, neuter: \emph{thermopōlium, thermopōliī n.} A word of comedy; such bars, with jars set in the counter, line the streets of excavated towns.

\begin{cousinweb}
\emph{thermopōlium} comes from Greek, built from \emph{thermos} 'hot' and \emph{pōleō} 'to sell' (Lewis \& Short give the sense, a place where warm drinks were sold). EtymDB does not list it.
\end{cousinweb}

\begin{grammarlens}[title={What you've built}]
One more word for telling, explaining and describing at B1.
\end{grammarlens}

\subsection*{Your turn}
\begin{itemize}
\raggedright
  \item \emph{Say it:} \emph{thermopōlium, thermopōliī}
  \item \emph{Say it:} \emph{thermopōlium, thermopōliī}, once more
  \item \emph{Say it:} a sentence of your own with \emph{thermopōlium, thermopōliī}
  \item \emph{Recall:} say the Latin for to go back and forth, then the Latin for near the city, then say \emph{thermopōlium, thermopōliī} again
\end{itemize}

\begin{tcolorbox}[breakable,colback=teal!4,colframe=teal!35!black,title={Before you move on}]
What does thermopōlium, thermopōliī mean? (\textquotedblleft{}A hot-drinks bar, a tavern selling warm drinks\textquotedblright{}.) Say it once more.
\end{tcolorbox}
\section[calda, caldae]{calda, caldae — hot water (for mixing with wine)}

\label{lesson:LA-C461-calda}



\begin{quote}
Before the new one: say the Latin for near the city, then the Latin for a hot-drinks bar.
\end{quote}

\subsection*{You'll want to know: calda, caldae}
\textbf{calda, caldae} — \textquotedblleft{}hot water (for mixing with wine)\textquotedblright{}.

\begin{quote}
Hospes caldam et vīnum rogat.
\end{quote}

\begin{quote}
\emph{The guest asks for hot water and wine.}
\end{quote}

First declension, feminine: \emph{calda, caldae f.}, short for \emph{(aqua) calida}. Romans liked wine mixed with hot water in winter.

\begin{cousinweb}
\emph{calda} is a shortened form of \emph{(aqua) calida} 'hot water', from \emph{calidus} 'hot' (Lewis \& Short), which is built on \emph{caleō} 'to be warm'.
\end{cousinweb}

\begin{grammarlens}[title={What you've built}]
One more word for telling, explaining and describing at B1.
\end{grammarlens}

\subsection*{Your turn}
\begin{itemize}
\raggedright
  \item \emph{Say it:} \emph{calda, caldae}
  \item \emph{Say it:} \emph{calda, caldae}, once more
  \item \emph{Say it:} a sentence of your own with \emph{calda, caldae}
  \item \emph{Recall:} say the Latin for near the city, then the Latin for a hot-drinks bar, then say \emph{calda, caldae} again
\end{itemize}

\begin{tcolorbox}[breakable,colback=teal!4,colframe=teal!35!black,title={Before you move on}]
What does calda, caldae mean? (\textquotedblleft{}Hot water (for mixing with wine)\textquotedblright{}.) Say it once more.
\end{tcolorbox}
\section[pōtō, pōtāre]{pōtō, pōtāre — to drink, to drink deep}

\label{lesson:LA-C461-poto}



\begin{quote}
Before the new one: say the Latin for a hot-drinks bar, then the Latin for hot water.
\end{quote}

\subsection*{You'll want to know: pōtō, pōtāre}
\textbf{pōtō, pōtāre} — \textquotedblleft{}to drink, to drink deep\textquotedblright{}.

\begin{quote}
Pōtōrēs usque ad mediam noctem pōtant.
\end{quote}

\begin{quote}
\emph{The drinkers drink on until midnight.}
\end{quote}

\emph{pōtō, pōtāre, pōtāvī, pōtātum} or \emph{pōtum}. Stronger than \emph{bibō}: drinking hard, drinking for pleasure. English \emph{potion} and \emph{potable} come from the same root.

\begin{cousinweb}
\emph{pōtō} is inherited from PIE \emph{peh\textsubscript{3}-} 'to drink', with Greek \emph{pinō} 'to drink' as a cousin (EtymDB; Lewis \& Short). English \emph{potion} and \emph{potable} come from it.
\end{cousinweb}

\begin{grammarlens}[title={What you've built}]
One more word for telling, explaining and describing at B1.
\end{grammarlens}

\subsection*{Your turn}
\begin{itemize}
\raggedright
  \item \emph{Say it:} \emph{pōtō, pōtāre}
  \item \emph{Say it:} \emph{pōtō, pōtāre}, once more
  \item \emph{Say it:} a sentence of your own with \emph{pōtō, pōtāre}
  \item \emph{Recall:} say the Latin for a hot-drinks bar, then the Latin for hot water, then say \emph{pōtō, pōtāre} again
\end{itemize}

\begin{tcolorbox}[breakable,colback=teal!4,colframe=teal!35!black,title={Before you move on}]
What does pōtō, pōtāre mean? (\textquotedblleft{}To drink, to drink deep\textquotedblright{}.) Say it once more.
\end{tcolorbox}
\section[fūmōsus, fūmōsa, fūmōsum]{fūmōsus, fūmōsa, fūmōsum — smoky, full of smoke}

\label{lesson:LA-C461-fumosus}



\begin{quote}
Before the new one: say the Latin for hot water, then the Latin for to drink.
\end{quote}

\subsection*{You'll want to know: fūmōsus, fūmōsa, fūmōsum}
\textbf{fūmōsus, fūmōsa, fūmōsum} — \textquotedblleft{}smoky, full of smoke\textquotedblright{}.

\begin{quote}
Popīna fūmōsa est, et oculī hospitum dolent.
\end{quote}

\begin{quote}
\emph{The cookshop is smoky, and the customers' eyes hurt.}
\end{quote}

An adjective of three endings, from \emph{fūmus} 'smoke' with \emph{-ōsus} 'full of'. Poets mock the smoky tavern.

\begin{cousinweb}
\emph{fūmōsus} is built on \emph{fūmus} 'smoke' with the suffix \emph{-ōsus} 'full of' (Lewis \& Short). English \emph{fume} comes from \emph{fūmus}.
\end{cousinweb}

\begin{grammarlens}[title={What you've built}]
One more word for telling, explaining and describing at B1.
\end{grammarlens}

\subsection*{Your turn}
\begin{itemize}
\raggedright
  \item \emph{Say it:} \emph{fūmōsus, fūmōsa, fūmōsum}
  \item \emph{Say it:} \emph{fūmōsus, fūmōsa, fūmōsum}, once more
  \item \emph{Say it:} a sentence of your own with \emph{fūmōsus, fūmōsa, fūmōsum}
  \item \emph{Recall:} say the Latin for hot water, then the Latin for to drink, then say \emph{fūmōsus, fūmōsa, fūmōsum} again
\end{itemize}

\begin{tcolorbox}[breakable,colback=teal!4,colframe=teal!35!black,title={Before you move on}]
What does fūmōsus, fūmōsa, fūmōsum mean? (\textquotedblleft{}Smoky, full of smoke\textquotedblright{}.) Say it once more.
\end{tcolorbox}
\section[immundus, immunda, immundum]{immundus, immunda, immundum — filthy, dirty, unclean}

\label{lesson:LA-C461-immundus}



\begin{quote}
Before the new one: say the Latin for to drink, then the Latin for smoky.
\end{quote}

\subsection*{You'll want to know: immundus, immunda, immundum}
\textbf{immundus, immunda, immundum} — \textquotedblleft{}filthy, dirty, unclean\textquotedblright{}.

\begin{quote}
Mēnsae immundae sunt, et cibus malus est.
\end{quote}

\begin{quote}
\emph{The tables are filthy, and the food is bad.}
\end{quote}

An adjective of three endings, from \emph{in-} 'not' and \emph{mundus} 'clean, neat'.

\begin{cousinweb}
\emph{immundus} is built from \emph{in-} 'not' and the adjective \emph{mundus} 'clean, neat' (EtymDB; Lewis \& Short).
\end{cousinweb}

\begin{grammarlens}[title={What you've built}]
One more word for telling, explaining and describing at B1.
\end{grammarlens}

\subsection*{Your turn}
\begin{itemize}
\raggedright
  \item \emph{Say it:} \emph{immundus, immunda, immundum}
  \item \emph{Say it:} \emph{immundus, immunda, immundum}, once more
  \item \emph{Say it:} a sentence of your own with \emph{immundus, immunda, immundum}
  \item \emph{Recall:} say the Latin for to drink, then the Latin for smoky, then say \emph{immundus, immunda, immundum} again
\end{itemize}

\begin{tcolorbox}[breakable,colback=teal!4,colframe=teal!35!black,title={Before you move on}]
What does immundus, immunda, immundum mean? (\textquotedblleft{}Filthy, dirty, unclean\textquotedblright{}.) Say it once more.
\end{tcolorbox}
